\documentclass[12pt, a4paper]{exam}
\usepackage{graphicx}
%\usepackage[left=0.8in, top=0.7in, total={6.2in,8in}]{geometry}
\usepackage[normalem]{ulem}
\renewcommand\ULthickness{1.0pt}   %%---> For changing thickness of underline
\setlength\ULdepth{1.3ex}%\maxdimen ---> For changing depth of underline
\usepackage{amssymb} 

\begin{document}
	%\thispagestyle{empty}
	\noindent
	\begin{minipage}[l]{0.1\textwidth}
		\noindent
		%\includegraphics[width=2.4\textwidth]{uct.png}
	\end{minipage}
\hfill
\begin{minipage}[c]{1.0\textwidth} % use 0.8 when using the image
	\begin{center}
		
		{
		\large \ \ \par
		\large \ \ \par
		\large  Alexander Cristaudo \par
		\small	CRSALE010	\par
		\large	MAM1019H	\par
	\large \textbf{Hand-in 6}	\par
\small	Date: September 21, 2022}
	\end{center}
\end{minipage}
\par
\vspace{0.2in}
\noindent
\uline{ \ \ \ \ \ \ \ \ \ \ \ \ \ \ \ \ \ \ \ \ \ \ \ \ \ \ \ \ \ \ \ \ \ \ \ \ \ \ \ \ \ \ \ \ \ \ \ \ \ \ \ \ \ \ \ \ \ \ \ \ \ \ \ \ \ \ \ \ \ \ \ \ \ \ \ \ \ \  \ \ \ \ \ \ \ \ \ \ \ \ \ \  \ \ \ \ \ \ \ \ \ \ \ \ \ \ \ \ \ \ \ \ \ \ }
\par 
\vspace{0.15in}
\noindent
\centering
{\small \bfseries 	Answers}

\begin{questions}
	\pointsdroppedatright
	\question
	\begin{parts}
		\part[9] False 
		\part[9] False
		\part False
		\part False
	\end{parts}
\vspace{0.2in}
	\pointsdroppedatright
	\question \ \\
	

\begin{tabular}{llll}
$n \cdot 3$ & $= n\cdot s(2)$ & \ \ \ \ \ \ \ \ \ \ \ \ \ \ \ \ & $s(2)=3$ \\
\ & $= n\cdot 2 + n$ & \ \ \ \ \ \ \ \ \ \ \ \ \ \ \ \ & (M2) with $m=2$ \\
\ & $= n\cdot s(1) + n$ & \ \ \ \ \ \ \ \ \ \ \ \ \ \ \ \ & $s(1)=2$ \\
\ & $= (n\cdot 1+n) + n$ & \ \ \ \ \ \ \ \ \ \ \ \ \ \ \ \ & (M2) with $m=1$ \\
\ & $= (n + n) + n$ & \ \ \ \ \ \ \ \ \ \ \ \ \ \ \ \ & We are given that $n\cdot 1 = n$  \\
\end{tabular} \\
This means $n \cdot 3 = (n+n)+n$


		\question \ \\
	\begin{tabular}{llll}
$2 \cdot 1$ & $= 2\cdot s(0)$ & \ \ \ \ \ \ \ \ \ \ \ \ \ \ \ \ & $s(0)=1$ \\
\ & $= 2\cdot 0 + 2$ & \ \ \ \ \ \ \ \ \ \ \ \ \ \ \ \ & (M2) with $n=2, m=0$ \\
\ & $= 0 + 2$ & \ \ \ \ \ \ \ \ \ \ \ \ \ \ \ \ & $2\cdot 0 = 0$ by (M1) \\
\ & $= 0 + s(1)$ & \ \ \ \ \ \ \ \ \ \ \ \ \ \ \ \ & $s(1)=2 $\\
\ & $= s(0+1)$ & \ \ \ \ \ \ \ \ \ \ \ \ \ \ \ \ & (A2) with $n=0, m=1$ \\
\ & $= s(0+s(0))$ & \ \ \ \ \ \ \ \ \ \ \ \ \ \ \ \ & $s(0)=1$ \\
\ & $= s(s(0+0))$ & \ \ \ \ \ \ \ \ \ \ \ \ \ \ \ \ & (A2) with $m=0, n=0$ \\
\ & $= s(s(0))$ & \ \ \ \ \ \ \ \ \ \ \ \ \ \ \ \ & $0+0=0$ by (A1) \\
\ & $= s(1)$ & \ \ \ \ \ \ \ \ \ \ \ \ \ \ \ \ & $s(0)=1$ \\
\ & $= 2$ & \ \ \ \ \ \ \ \ \ \ \ \ \ \ \ \ & $s(1)=2$ \\


\end{tabular} \\	
This means that $2\cdot 1 = 2$		
		
\end{questions}
\vspace{0.75in}
{\large \bfseries }
\end{document}